% ===== PRÉAMBULE CANONIQUE — à coller VERBATIM en tête de CHAQUE .tex =====
\documentclass[11pt,a4paper]{article}
\usepackage[utf8]{inputenc}\usepackage[T1]{fontenc}\usepackage{lmodern}
\usepackage[french]{babel}
\usepackage{amsmath,amssymb,mathtools}\usepackage{icomma}
\usepackage{siunitx}\sisetup{locale=FR}  % \qty{}{}, \si{}, \ang{}
\usepackage{xcolor}\usepackage{colortbl}
\usepackage{tikz}\usepackage{pgfplots}\pgfplotsset{compat=1.18}
\usepackage[siunitx,europeanresistors]{circuitikz} % circuits (élec.)
\usepackage[most]{tcolorbox}\usepackage{enumitem}\usepackage{array,booktabs}
\usepackage[a4paper,margin=2.2cm]{geometry}
\usetikzlibrary{arrows.meta,calc,angles,quotes,positioning,patterns,%
  backgrounds,shapes.geometric,decorations.pathmorphing,decorations.markings,intersections}
\definecolor{primary}{RGB}{222,49,99}\definecolor{primarydark}{RGB}{150,22,55}
\definecolor{defcol}{RGB}{0,114,178}\definecolor{methcol}{RGB}{52,92,148}
\definecolor{excol}{RGB}{0,135,90}\definecolor{attncol}{RGB}{200,40,55}
\tcbset{boxbase/.style={enhanced,breakable,boxrule=0.9pt,arc=2.5pt,
  left=3mm,right=3mm,top=2.3mm,bottom=2.3mm,fonttitle=\bfseries,coltitle=white}}
% --- boîtes TUTORIEL (noms EXACTS, ne pas renommer) ---
\newtcolorbox{cle}[1][]{boxbase,colback=defcol!6,colframe=defcol,
  title={\textsc{À retenir}\ifx\relax#1\relax\relax\else\textmd{ : \itshape #1}\fi}}
\newtcolorbox{methode}[1][]{boxbase,colback=methcol!7,colframe=methcol,sharp corners,
  title={\textsc{Méthode}\ifx\relax#1\relax\relax\else\textmd{ --- \itshape #1}\fi}}
\newtcolorbox{exemplebox}[1][]{boxbase,colback=excol!5,colframe=excol,
  title={\textsc{Exemple}\ifx\relax#1\relax\relax\else\textmd{ : \itshape #1}\fi}}
\newtcolorbox{piege}[1][]{boxbase,colback=attncol!6,colframe=attncol,
  title={\textsc{Piège}\ifx\relax#1\relax\relax\else\textmd{ : \itshape #1}\fi}}
\newtcolorbox{remarque}[1][]{boxbase,colback=black!4,colframe=black!45,
  title={\textsc{Remarque}\ifx\relax#1\relax\relax\else\textmd{ : \itshape #1}\fi}}
% --- boîtes EXERCICE / CORRIGÉ (noms EXACTS, auto-comptées) ---
\newtcolorbox[auto counter]{exo}[1][]{boxbase,colback=primary!4,colframe=primary,
  title={\textsc{Exercice}~\thetcbcounter\ifx\relax#1\relax\relax\else\textmd{ --- \itshape #1}\fi}}
\newtcolorbox[auto counter]{corrige}[1][]{boxbase,colback=excol!4,colframe=excol,
  title={\textsc{Corrigé}~\thetcbcounter\ifx\relax#1\relax\relax\else\textmd{ --- \itshape #1}\fi}}
\newcommand{\vv}[1]{\vec{#1}}\newcommand{\ds}{\displaystyle}
\newcommand{\R}{\mathbb{R}}\newcommand{\N}{\mathbb{N}}\newcommand{\Z}{\mathbb{Z}}
% =========================================================================
\begin{document}
\sisetup{per-mode=symbol}
% ================= DIFFICILE =================

\begin{exo}[l'énergie d'un satellite --- vrai ou faux]
Un satellite est en orbite circulaire autour de la Terre. Indique si chaque
affirmation est \textbf{vraie} ou \textbf{fausse}.
\begin{enumerate}[label=\alph*)]
  \item Le satellite possède une énergie cinétique.
  \item Sur son orbite circulaire, son énergie cinétique $E_c$ reste constante.
  \item La somme $E_c+E_p$ est la même en tout point de l'orbite.
  \item L'énergie mécanique $E_m$ du satellite est positive.
\end{enumerate}
\end{exo}

\begin{exo}[bilan d'un satellite géostationnaire]
Un satellite géostationnaire a une masse $m=\qty{1000}{\kilogram}$ et gravite à
$r=\qty{4.22e7}{\metre}$ ($\mathcal{G}M_T=\qty{4.0e14}{}$ SI, donc
$\mathcal{G}M_T m=\qty{4.0e17}{}$).
\begin{enumerate}[label=\alph*)]
  \item Calcule son énergie cinétique $E_c$.
  \item Calcule son énergie potentielle $E_p$.
  \item Calcule $E_m$ et vérifie la relation $E_m=-E_c$.
\end{enumerate}
\end{exo}

\begin{exo}[libérer un satellite]
Un satellite de masse $m=\qty{500}{\kilogram}$ orbite à $r=\qty{1.0e7}{\metre}$
($\mathcal{G}M_T=\qty{4.0e14}{}$ SI).
\begin{enumerate}[label=\alph*)]
  \item Calcule son énergie mécanique $E_m=-\dfrac{\mathcal{G}M_T m}{2r}$.
  \item Quelle énergie faut-il lui fournir pour l'envoyer « à l'infini »
        (c'est-à-dire amener $E_m$ à $0$) ?
\end{enumerate}
\end{exo}

\begin{exo}[lequel est le plus fortement lié ?]
Deux satellites de \emph{même} masse orbitent : $A$ à $r_A=\qty{2.0e7}{\metre}$ et
$B$ à $r_B=\qty{1.0e7}{\metre}$.
\begin{enumerate}[label=\alph*)]
  \item À partir de $E_m=-\dfrac{\mathcal{G}M_T m}{2r}$, établis le rapport
        $\dfrac{E_{mA}}{E_{mB}}$.
  \item Calcule-le, puis dis lequel a l'énergie mécanique la plus \emph{négative}.
  \item Lequel est le plus fortement lié à la Terre (le plus difficile à libérer) ?
\end{enumerate}
\end{exo}

\end{document}
